% defined-in: zorich (page 398)
% entity-id: thm-cauchy-criterion-vector-valued
% entity-type: theorem

\hypertarget{thm-cauchy-criterion-vector-valued}{}
\hypertarget{thm-cauchy-criterion-vector-valued}{}
\begin{theorem}[Cauchy Criterion for Vector-Valued Functions]
\[
\mForall{X} \mForall{\mathcal{B} \subseteq \mathcal{P}(X)} \mForall{\entityref{obj-function}{f} \colon X \mTo \entityref{obj-euclidean-space}{\mathbb{R}^n}}
\left( \mExists{\entityref{op-limit}{\lim}_{\mathcal{B}} f(x)} \mIff \mForall{\varepsilon > 0} \mExists{B \in \mathcal{B}} \left( \omega(\entityref{obj-function}{f}; B) < \varepsilon \right) \right)
\]
\[
\text{where } \omega(\entityref{obj-function}{f}; B) \mDefinedAs \entityref{op-supremum}{\sup}_{x_1, x_2 \in B} \entityref{op-euclidean-distance}{d}(f(x_1), f(x_2))
\]
\end{theorem}

\begin{proof}[RU]
Доказательство опирается на полноту пространства $\entityref{obj-euclidean-space}{\mathbb{R}^n}$. Пусть $\entityref{obj-euclidean-distance}{d}(f(x_1), f(x_2))$ обозначает евклидово расстояние. Условие $\mForall{\varepsilon > 0} \mExists{B \in \mathcal{B}} (\omega(\entityref{obj-function}{f}; B) < \varepsilon)$ означает, что \entityref{obj-image}{образ} \entityref{obj-set}{множества} $B$ при \entityref{obj-function}{отображении} $\entityref{obj-function}{f}$ становится сколь угодно малым в смысле диаметра. Поскольку $\entityref{obj-euclidean-space}{\mathbb{R}^n}$ является полным метрическим пространством, выполнение критерия Коши для базы $\mathcal{B}$ необходимо и достаточно для существования предела \entityref{obj-function}{функции} $\entityref{obj-function}{f}$ по этой базе. Доказательство аналогично одномерному случаю, где вместо модуля разности используется $\entityref{op-euclidean-distance}{d}(f(x_1), f(x_2))$.
\end{proof}

\begin{proof}[EN]
The proof relies on the completeness of the space $\entityref{obj-euclidean-space}{\mathbb{R}^n}$. Let $\entityref{op-euclidean-distance}{d}(f(x_1), f(x_2))$ denote the Euclidean distance. The condition $\mForall{\varepsilon > 0} \mExists{B \in \mathcal{B}} (\omega(\entityref{obj-function}{f}; B) < \varepsilon)$ implies that the image of the set $B$ under the mapping $\entityref{obj-function}{f}$ becomes arbitrarily small in terms of its diameter. Since $\entityref{obj-euclidean-space}{\mathbb{R}^n}$ is a complete metric space, the Cauchy criterion for the base $\mathcal{B}$ is necessary and sufficient for the existence of the limit of the function $\entityref{obj-function}{f}$ along this base. The proof is analogous to the one-dimensional case, replacing the absolute difference with the distance $\entityref{op-euclidean-distance}{d}(f(x_1), f(x_2))$.
\end{proof}